\section{Ethics Considerations}
\label{sec:ethics-considerations}
% This work does not raise any ethical concern.

%%
This work is a Systematization of Knowledge (SoK) on web tracking. 
%
As such, it does not involve any data collection, user studies, or measurement and deployment of tracking technologies.
%
Instead, it systematizes developments in the field on online tracking through research and industrial advancements to identify gaps and open problems.
%
Nevertheless, we discuss the ethical implications of our work in shaping the future research directions and influencing the involved stakeholders.
%
We analyze the relevant stakeholders, assess the potential positive and negative impacts of this research, and justify why conducting and publishing this SoK is ethically warranted.\\

\noindent \textbf{Stakeholder Analysis.} 
%
Our work discusses six stakeholders: 
%
(1) \textit{Online users}: 
%
are individuals whose web activity may be observed, inferred, or profiled using different tracking mechanisms. 
%
(2) \textit{Browser vendors}: 
%
are the developers of browsers that mediate access to the web, while designing and deploying technical countermeasures to balance privacy, usability, performance, and compatibility.
%
(3) \textit{Website publishers}: 
%
operate websites and monetize their content by relying on analytics, personalization, and advertising, while maintaining functionality under evolving privacy constraints.
%
(4) \textit{Advertising and tracking platforms}: 
%
that design and operate tracking infrastructures and are affected by regulatory changes, browser policies, and public scrutiny.
%
(5) \textit{Regulators and policymakers}: 
%
are responsible for enforcing data protection and competition laws. They rely on accurate technical understanding to craft effective and enforceable regulation.
%
(6) \textit{Research community}: 
%
comprises of academic and industrial researchers studying web tracking and privacy to understand, measure, and defend against evolving tracking techniques.\\

%%
\noindent \textbf{Impact of the Research.} 
%
Next, we list the positive and negative impacts of our research on different stakeholders.

%%
\noindent \textbf{\textit{Positive Impacts.}} 
%
(1) \textit{Improved Understanding (Impact on Users, Regulators, and Researchers)}: 
%
By systematizing advances in tracking techniques, defenses, and measurement methodologies, this work benefits users and regulators by clarifying what forms of tracking are known, which ones are mitigated, and which ones remain unresolved. 
%
This SoK also serves as a unified source of knowledge for early career researchers, users, and regulators to understand technicalities of different tracking mechanisms. 
%
(2) \textit{Future Research Directions (Impact on Researchers and Browser Vendors)}: 
%
Our work identifies underexplored areas and open problems based on advancements in tracking landscape allowing new as well as established researchers to focus on important directions of research to improve the state of web tracking rather than making marginal improvements or working well-explored problems. 
%
Similarly, it can also inform browser vendors to anticipate emerging forms of web tracking and work towards in-browser countermeasures. 
%
(3) \textit{Informing Platform and Policy Design (Impact on Browser Vendors and Regulators)}: 
%
By synthesizing years of developments in web tracking, this SoK informs browser vendors and policymakers about unintended consequences and gaps in current approaches, while guiding future policy changes---both browser-level as well as legal. 

%%
\noindent \textbf{\textit{Negative Impacts.}} 
%
(1) \textit{Potential for Adversarial Use (Impact on Users)}: 
%
This paper can inform motivated adversaries (i.e., trackers) to exploit the identified gaps in online tracking research by deploying them within their tracking technologies to circumvent existing protections. 
%
Identified open problems may further motivate to make these mechanisms stealthier.
%
(2) \textit{Misinterpretation or Selective Use of Findings (Impact on Policy)}: 
%
Industry or policy actors may selectively cite parts of this SoK to justify particular technical or regulatory choices, potentially oversimplifying nuanced trade-offs discussed in the specific academic work in the literature.\\ 


%%
\noindent \textbf{Justification for Research.}
%
Despite these risks, we argue that publishing this SoK is ethically justified and necessary. 
%
Web tracking is already widely deployed at scale, often with limited transparency to users and regulators. 
%
The primary ethical risk lies not in documenting this ecosystem, but in allowing knowledge fragmentation or obscurity to persist. 
%
Several factors mitigate the highlighted potential negative impacts. 
%
First, this work does not introduce new tracking mechanisms, datasets, or attack surfaces; it systematizes and critically evaluates information that is already publicly available on the internet or published in the research literature. 
%
Second, by explicitly framing open problems, we aim to encourage research that prioritizes user protection, accountability, and robustness rather than exploitation. 
%
Finally, increased visibility into the limitations of existing defenses and measurements enables more informed decision-making by browser vendors, regulators, and researchers, reducing the likelihood of misleading claims about privacy guarantees. 
%
Overall, this SoK contributes to an ethically grounded understanding of web tracking. 
%
By establishing what is known, what remains unsolved, and where future research efforts are needed the most, this work supports long-term improvements in web privacy that outweigh the potential risks from systematizing this knowledge.